% ###################
% ## Filename: sCM_inventory.tex
% ## Location: D:\Project\Prej_Teaching\ISE\SupplyChainMgt\presentation\scm_02_inventory
% ## Created: Jan 26, 2011
% ## Last updated: Feb 2, 2010
% ## Figure Location: ..\Figures\
% ## Note: Presentation, Teaching
% ###################

%\documentclass[10pt,presentation]{beamer}
%\documentclass[10pt,handout]{beamer}
\documentclass[10pt,handout,notes=show]{beamer}

%============================================================================================================--
\usepackage{fontspec}
\usepackage{xunicode}
\usepackage{xltxtra}
\usepackage{array}
\usepackage{xcolor}
\usepackage{verbatim}
\usepackage{beamerthemesplit}
\usepackage{beamerbasetemplates}
\usepackage{graphics}
\usepackage{graphicx}
\usepackage{hyperref}
\usepackage{multirow}
\usepackage{tikz,pgfplots}
\usepackage{algorithm}
\usepackage{algorithmic}
\usepackage{soul,xcolor}
\usepackage{colortbl}
\usepackage{multicol}

\XeTeXlinebreaklocale"th"
\defaultfontfeatures{Scale=1.5, Mapping=tex-text} %{Scale=1.23}



\usetheme{Madrid}  % JuanLesPins  Rochester
           %  Berkeley  Palo Alto    with sidebar and top
           % Goettingen  Marburg     with sidebar
           %  Copenhagen  Luebeck  Warsaw

\setbeamertemplate{headline}
{%
    \begin{beamercolorbox}{section in head/foot}
        \vskip-2pt \hskip5pt
        \insertsectionnavigationhorizontal{\textwidth}{}{}
        \hskip5pt \vskip-2pt
    \end{beamercolorbox}
}%
\setbeamertemplate{footline}
{%
    \begin{beamercolorbox}{section in head/foot}
        \vskip-2pt \hskip5pt
        \insertshorttitle \hfill \insertshortdate \hskip10pt \insertframenumber / \inserttotalframenumber
        \vskip-2pt\hskip5pt
    \end{beamercolorbox}
}%
\definecolor{alertedOrange}{RGB}{255,36,0}
\setbeamercolor{alerted text}{fg=alertedOrange}
\useinnertheme{rounded}
\usecolortheme{rose}
%\setbeamercovered{transparent}
\usefonttheme{structuresmallcapsserif}
\definecolor{LRed}{rgb} {1,.8,.8}
\definecolor{LBlue}{rgb}{.8,.8,1}

%%%%%%%%%%%%%%%%%%%%%%%%%%%%%%%%%%%%%%%%%%%%%%%%%%%%%%%%%%%%%%%%%%%%%%%%%%%%%%%
% Notation:
\newtheorem{observation}{Observation}
\newcommand{\figureLocation}[1]{../figures/#1}
\newcommand{\figureLoc}[1]{../figures/#1}
\newcommand{\refLocation}[1]{../../../../Reference/#1}
\newcommand{\NextLine}{\vskip 0.05\textheight}
\newcommand{\NextHalfLine}{\vskip 0.02\textheight}
\newcommand{\RemoveLine}{\vskip -0.1\textheight}
\newcommand{\RemoveHalfLine}{\vskip -0.02\textheight}
\newcommand{\BackTab}{\hskip -0.15\textwidth}
\newcommand{\HalfBackTab}{\hskip -0.11\textwidth}
\newcommand{\DualFigureSize}{1.10\textwidth}
\newcommand{\MiniPageSize}{0.52\textwidth}
\newcommand{\code}[1]{\textcolor{blue}{#1} }
\newcommand{\rightNote}[2]{\hfill \uncover<#1>{\textbf{\code{\#\# {#2}  }}}}
\newcolumntype{P}[1]{>{\raggedright\arraybackslash}p{#1}<{}}

\author[Kittithreerapronchai]{%
  Oran~Kittithreerapronchai\inst{1} }
\institute[IE Chula]{
  \inst{1}%
  Department of Industrial Engineering, Chulalongkorn University\\
  Bangkok 10330 THAILAND \\
  }
\subject{}
\date{}

\setbeamertemplate{enumerate items}[default]
\setbeamertemplate{bibliography item}[text]
\bibliographystyle{alpha}
\newcommand{\textFont}[1]{\footnotesize{#1}}
\setcounter{MaxMatrixCols}{20}

%%%--------------------------------

\makeatletter
\defbeamertemplate{note page}{infolines}
{%
  {%
    \small
    \insertvrule{.25\paperheight}{white!90!black}
    \vskip-.25\paperheight
    \nointerlineskip
    \vbox{
      \hfill\insertslideintonotes{0.25}\hskip-\Gm@rmargin\hskip0pt%
      \vskip-0.25\paperheight%
      \nointerlineskip
      \begin{pgfpicture}{0cm}{0cm}{0cm}{0cm}
        \begin{pgflowlevelscope}{\pgftransformrotate{90}}
          {\pgftransformshift{\pgfpoint{-2cm}{0.2cm}}%
          \pgftext[base,left]{\footnotesize\the\year-\ifnum\month<10\relax0\fi\the\month-\ifnum\day<10\relax0\fi\the\day}}
        \end{pgflowlevelscope}
      \end{pgfpicture}}
    \nointerlineskip
    \vbox to .25\paperheight{\vskip0.5em
      \hbox{\insertshorttitle[width=8cm]}%
      \setbox\beamer@tempbox=\hbox{\insertsection}%
      \hbox{\ifdim\wd\beamer@tempbox>1pt{\hskip4pt\raise3pt\hbox{\vrule
            width0.4pt height7pt\vrule width 9pt
            height0.4pt}}\hskip1pt\hbox{\begin{minipage}[t]{7.5cm}\def\breakhere{}\insertsection\end{minipage}}\fi%
      }%
      \setbox\beamer@tempbox=\hbox{\insertsubsection}%
      \hbox{\ifdim\wd\beamer@tempbox>1pt{\hskip17.4pt\raise3pt\hbox{\vrule
            width0.4pt height7pt\vrule width 9pt
            height0.4pt}}\hskip1pt\hbox{\begin{minipage}[t]{7.5cm}\def\breakhere{}\insertsubsection\end{minipage}}\fi%
      }%
      \setbox\beamer@tempbox=\hbox{\insertshortframetitle}%
      \hbox{\ifdim\wd\beamer@tempbox>1pt{\hskip30.8pt\raise3pt\hbox{\vrule
            width0.4pt height7pt\vrule width 9pt
            height0.4pt}}\hskip1pt\hbox{\insertshortframetitle[width=7cm]}\fi%
      }%
      \vfil}%
  }%
  \vskip.25em
  \nointerlineskip
  %\setmainfont{TH Sarabun New}
  \begin{minipage}{\textwidth} % this is an addition
  \footnotesize{\fontspec{TH Sarabun New} \insertnote}
  \end{minipage}               % this is an addition
  %\setmainfont{Computer Modern}

}
\makeatother

\setbeamertemplate{note page}[infolines]

\title[LSCM v3.0]{Lecture 06:\\ Inventory Management \\ ~}
\hypersetup{
  pdftitle={06 Inventory Management},
  pdfsubject={2104509lscm & 2184408scm handout},
  pdfproducer={XeLatex/WinEdt},
  pdfcreator={oran.k@chula.ac.th}
}


\begin{document}

%============================================================================================================--
\begin{frame}[t]
  \titlepage \label{Slide:Title}

  \NextLine   \NextLine
  \begin{flushright}
    \tiny{ \textbf{last updated:} \today } \hskip3ex
  \end{flushright}

  \note{
    \noindent\textbf{Purpose:} Introduce to inventory control\par
    \noindent\textbf{Leading Question:} Where have you encountered this kind of supply-chain decision?\par
    \noindent\textbf{Teaching cue:} Get one example and identify who owns the decision and what outcome they want.\par
    \begin{itemize}
      \item $\exists$ Inventory Analysis class $\rightarrow$ theory
      \item \textbf{Scope} review key concept and highlight conflict areas between academic and practice
    \end{itemize}
  }
\end{frame}



%============================================================================================================--
\begin{frame}[t]
  \frametitle{Outline}

  \tableofcontents
  \vfill
  \begin{flushright}
    \tiny{ \textbf{Key Ref.:}~\cite{jacobs2010operations}~\cite{ballou2007evolution} \cite{chopra2007supply} \cite{goetschalckx2011supply}} \hskip3ex
  \end{flushright}


  \note{
    \noindent\textbf{Purpose:} Overview Inventory in SCM\par
    \noindent\textbf{Leading Question:} Which topic would you consult first for the decision raised at the opening?\par
    \noindent\textbf{Teaching cue:} Point to the relevant section and ask which question should be resolved by the end.\par
    \begin{itemize}
      \item \textbf{Inventory for?} key inventory concept that lead to division
      \item \textbf{EOQ:} The most important 70yr-old trade-off model in IE for certainty demand
      \item \textbf{Review System:} how practically manage inventory
      \item \textbf{Newsvendor Model:} inventory for perishable items/ uncertainty demand
      \item important and motivation of SCM~\cite{jacobs2010operations}~\cite{ballou2007evolution} \cite{chopra2007supply} \cite{goetschalckx2011supply}

    \end{itemize}
  }
\end{frame}

%================================================================================================-
\section[Intro]{Why do we have an inventory}
%================================================================================================-
\begin{frame}
  \frametitle{What is Inventory?}
  \begin{description}
      \item[Inventory] is a the stock of any \alert{item} or \alert{resource used} in an organization.
      \item[] For example, raw materials, finished products, components, parts, supplies, and work-in-process (WIP)
  \end{description}

  \NextLine
  \textbf{Roles of Inventory}
  \begin{itemize}
      \item To prevent \alert{variability} in \alert{demands} and \alert{delivery time}
      \item To allow \alert{flexibility in production} scheduling
      \item To take advantage of \alert{economic purchase-order} size
  \end{itemize}

  %%%------
  \note{
    \noindent\textbf{Purpose:} Remind the definition\par
    \noindent\textbf{Leading Question:} What physical and financial consequences follow from holding a unit in stock?\par
    \noindent\textbf{Teaching cue:} Use a part on a shelf to trace ownership, working capital, and service protection.\par
    \begin{itemize}
      \item \textbf{Uncertainty}: demand; lead times; delivery of supply
      \item \textbf{Transportation and Logistics:} inventory in transit from one location to other
      \item \textbf{Speculation}: FG held in anticipation of a rise in their value or costs
      \item \textbf{Smoothing}: irregularities of demand pattern
      \item \textbf{Economy}: cheaper to produce/order in large batches; setup cost and time
    \end{itemize}
    \textbf{\textsc{Inventory as Managerial Tools}}
    \begin{itemize}
    \item Managing capacity
    \begin{itemize} \scriptsize{
      \item time flexibility from workforce (OT, pay per piece)
      \item seasonal workforce or sub-contracting
      \item designing product flexibility into production processes }
    \end{itemize}
    \item Managing inventory
    \begin{itemize} \scriptsize{
      \item common components across multiple products
      \item building inventory of high demand or predictable demand products }
    \end{itemize}
  \end{itemize}

  }
\end{frame}


%=================================================================================================
\begin{frame}
  \frametitle{Demand $\pm$ Supply = Inventory }


  \begin{center}
      \includegraphics [width=0.7\textwidth]{\figureLocation{scm_inv/demandSupply.png}}
  \end{center}


  \NextLine
  \textsc{Type of Inventory}
  \begin{itemize}
    \item \textbf{Pipeline inventory:} \alert{transportation}, minimal lot-sizing, e.g., vessel/train
    \item \textbf{Buffer inventory:}  imbalance of process, \alert{bottleneck}
    \item \textbf{Seasonal inventory:} growth/slow down, \alert{speculation}, e.g, BBQ grill, A/C
    \item \textbf{Cycle inventory:} regular consumption to exploit economy of scale
    \item \textbf{Safety inventory:} \alert{perspective} risks
  \end{itemize}


  \note{
    \noindent\textbf{Purpose:} Different name, different purpose of inventory\par
    \noindent\textbf{Leading Question:} Which terms belong in an inventory balance over one period?\par
    \noindent\textbf{Teaching cue:} Reconcile opening stock plus receipts minus demand served with closing stock.\par
    \begin{itemize}
      \item \textbf{Qty. Supply $\neq$ Qty. Demand:} growth/slow down, fluctuation, \alert{perspective} risks
      \item \textbf{Lead time:} \alert{transportation}, minimal lot-sizing
      \item \textbf{Others:} sale opportunity, expected shortage
    \end{itemize}
  }
\end{frame}

%================================================================================================-
\begin{frame}
  \frametitle{Why inventory is bad? }

  \begin{center}
      \includegraphics [width=0.6\textwidth]{\figureLocation{scm_inv/inventoryHidding.jpg}}
  \end{center}

 \begin{itemize}
   \item \textbf{Capital:} opportunity (i.e., raw materials, labor, money), obsolete
   \item \textbf{Physical space:} need shelter and system to take care
   \item \textbf{Liability:} risk of theft/fire \hfill \hyperlink{questions}{Pg:Inventory Question}
 \end{itemize}

  \note{
    \noindent\textbf{Purpose:} Key concepts of inventory\par
    \noindent\textbf{Leading Question:} When does inventory support service and when can it hide a problem?\par
    \noindent\textbf{Teaching cue:} Compare a purposeful buffer with excess stock caused by unstable replenishment.\par
    \begin{itemize}
      \item \textbf{capital:} money, time, resources
      \item \textbf{space} more inventory $\rightarrow$  difficult to manage, your own space, outsoucing $\rightarrow$ transportation
      \item \textbf{perishable}: short life cycle, burst-or-boom
      \item \textbf{liability} need to take care, control
      \item \textbf{inventory = evil?} right quantity, more $>$ less,
    \end{itemize}
  }
\end{frame}

%============================================================================================================
\usetikzlibrary{arrows,positioning,decorations.pathmorphing,calc}
\renewcommand{\textFont}[1]{\footnotesize{#1}}
\tikzstyle{lvBlock}    =[rectangle,draw=black!80,thick,text centered,text width=25mm,text height=5mm] %rounded  corners=3mm,
\tikzstyle{detailBlock}=[rectangle,draw=black!80,text width=70mm,text height=5mm]
%=================================================================================================
\begin{frame}[label=questions]
  \frametitle{Hierarchy of inventory decision \forIE}


  \begin{center}
    \begin{tikzpicture}[xscale=4,yscale=2]
      \node[]       (ref) {\textFont{~}};
      \node[lvBlock,fill=black!15] (lv2) [left  = of ref]  {\textbf{Deployment Decision}};
      \node[lvBlock,fill=red!25]   (lv1) [above = of lv2]  {\textbf{Supply Chain Decision}};
      \node[lvBlock,fill=blue!25]  (lv3) [below = of lv2]  {\textbf{Replenishment Decision}};
      \node[detailBlock,fill=red!25](detail1)  [right = of lv1]  {\textFont{
      \begin{tabular}{l}
        $\bullet$ What are the potential \alert{alternatives}? \\
        $\bullet$ How should the product be \alert{designed}?
      \end{tabular}
      }};
      \node[detailBlock,fill=black!15](detail2)  [right = of lv2]  {\textFont{
      \begin{tabular}{l}
        $\bullet$ What \alert{SKUs} should be carried as inventory?\\
        $\bullet$ In what \alert{form} should they be maintained?\\
        $\bullet$ How much of each should be held and where?
      \end{tabular}
      }};
      \node[detailBlock,fill=blue!25](detail3)  [right = of lv3] {\textFont{
      \begin{tabular}{l}
        $\bullet$ How often should inventory status be determined?\\
        $\bullet$ When should a replenishment decision be made?\\
        $\bullet$ What replenishment qty should be?
      \end{tabular}
      }};
      \draw[->,thick] (lv1) -- (lv2);
      \draw[->,thick] (lv2) -- (lv3);
    \end{tikzpicture}
  \end{center}


  \note{
    \noindent\textbf{Purpose:} Alternative before inventory storage\par
    \noindent\textbf{Leading Question:} Which inventory choices are made by item, location, and network?\par
    \noindent\textbf{Teaching cue:} Assign a service target, location choice, and order size to decision owners.\par
    \begin{itemize}
      \item \textbf{Str: Potent. Altern.:} what consumption? e.g., media $\rightarrow$ internet steaming
      \item \textbf{Str: Design. Product:} design to reduce component? family group, common VS module; decided by mrkt/designer
      \item \textbf{Deploy: Which items?} risk VS space \& money; decided by
      \item \textbf{Deploy: What form?} common for every model VS module for quick response
      \item \textbf{Deploy: Where/Max Qty} money VS responce ;
      \item \textbf{Replenish: Check/Qty/When:} common oper. decision; decided by
    \end{itemize}
  }
\end{frame}

%================================================================================================-
\section[EOQ]{Economic Order Quantity}
%================================================================================================-
\begin{frame}
  \frametitle{Main inventory control questions in IE}

  \textsc{Replenishment Decisions}
  \begin{itemize}
    \item[] \hskip3ex \code{Given,} what items need to stock with perfect information \uncover<2 |handout:1>{\hfill $\rightarrow$ \code{review system} }
    \item \alert{When to order} items from supplier or when to replenish?
    \item \alert{How much to order} from supplier on each ordering
  \end{itemize}

  \NextLine
  \textsc{Goal:}
  \begin{alertblock}{}{}
    Keep \alert{enough inventory} to meet customer demand and also \alert{cost effective}
  \end{alertblock}

  \NextLine
  \uncover<2->{
    \textsc{Two Extreme Cases: Independent VS Dependent}
    \begin{itemize}
      \item \textbf{Independent}: demands for various items are \alert{unrelated} to others
      \item \textbf{Dependent}: demands for any one item is a proportional to others
    \end{itemize}
  }

  \note{
    \noindent\textbf{Purpose:} Deploy replenishment scope\par
    \noindent\textbf{Leading Question:} What is the difference between how much to order and when to order?\par
    \noindent\textbf{Teaching cue:} Identify an order quantity and a reorder threshold for the same item.\par
    \begin{itemize}
      \item \textbf{What item:}  sku  $\rightarrow^{policy}$ impact \& high value; guild the next two question
      \item \textbf{When to order:} trigger point $\rightarrow^{tech}$ reviewing process (inventory or duration + other info)
      \item \textbf{How much:} qty $\rightarrow^{cost}$ cycle inventory + 'buffer'
      \item \textbf{Object:} a practical objective; inventory $\approx$ demand with opt cost
      \item \textbf{Why not dependence?} we can proportion --- difficult to assign cost
    \end{itemize}
  }
\end{frame}

%================================================================================================-
\begin{frame}
  \frametitle{Inventory costs}


  \begin{center}
    \includegraphics[width=0.45\linewidth]{\figureLocation{scm_inv/tradeOff.png}}\\
  \end{center}

 \begin{description}
   \item[Holding costs] costs for storage, capital, handling, spoilage, insurance
   \item[Ordering costs] costs of placing an order (fixed \& variable)
   \item[Shortage costs] costs of running out, back-order, lost-sale
 \end{description}

  \note{
    \noindent\textbf{Purpose:} Re-Cap/Overview about inventory cost/ difference between academic VS practical\par
    \noindent\textbf{Leading Question:} Which costs move when order frequency changes?\par
    \noindent\textbf{Teaching cue:} Classify purchasing, fixed order, shortage, and holding costs and specify their units.\par
    \begin{itemize}
      \item All of inventory cost (holding+ordering+backloging) is opportunity cost
      \item \textbf{Backloging} ($USD/time$ or $USD/unit-period$):  cost associated with missing shipment/consumption $\rightarrow$ most difficult to estimate
      \begin{itemize} \tiny{
        \item back-order case:- orders cannot be filled are held on the books until the next shipment arrives.
        \item lost-sale case:- orders cannot be filled are lost. }
      \end{itemize}
      \item \textbf{Holding} ($USD/unit-period$): cost associated with carrying inventory for a specified period of time $\rightarrow$ easiest to estimate
      \begin{itemize} \tiny{
        \item capital cost:- opportunity cost of alternative investment
  	    \item storage cost:- physical space to store the items
  	    \item risk cost:- breakage, spoilage, deterioration, and obsolescence}
      \end{itemize}
      \item \textbf{Setup/Ordering} ($USD/time$: costs associated with ordering process
      \item[] \hskip 3ex It do not vary with the size of the order but are incurred each time an order is placed, e.g., administrative cost, a trucking cost
      \begin{itemize} \tiny{
        \item fixed cost:- independent of lot size of the order as long as it  is not zero
        \item variable cost:- a per-unit basis}
      \end{itemize}
    \end{itemize}
  }
\end{frame}


%================================================================================================-
\begin{frame}
  \frametitle{Economic Order Quantity: Simplest model}

  \begin{itemize}
    \item single and independent product
    \item uniform (continuous) and constant demand
    \item instant replenishment and constant lead time
  \end{itemize}

  \begin{center}
    \includegraphics[width=0.85\linewidth]{\figureLocation{scm_inv/EOQ_SawTooth.jpg}}\\
  \end{center}
  \begin{flushright}
    \tiny{\textbf{source:} Chase and Jacob. 2011.\hskip10em}\\
  \end{flushright}

  \note{
    \noindent\textbf{Purpose:} Introduce to EOQ\par
    \noindent\textbf{Leading Question:} Under what assumptions is the economic order quantity defensible?\par
    \noindent\textbf{Teaching cue:} Check demand stability, fixed order cost, replenishment timing, and absence of constraints.\par
    \begin{itemize}
      \item \textbf{Single \& Independence:} no interference or correlation from other SKUs; pattern also repeat
      \item \textbf{Uniform \& Constant Demand:} same demand/period every period
      \item \textbf{Instant \& Constant LT:} no variability, no buffer no safety stock \\[8pt]
      \item \textbf{Decision Variable} is \alert{ordering quantity}; example monthly demand is 10 units or 120 units annually.
    \end{itemize}
  }
\end{frame}

%================================================================================================-
\begin{frame}
  \frametitle{Derivation of EOQ }

  \begin{itemize}
    \item \textbf{Parameters:}
    \begin{itemize}
      \item Annual demand ($D$)
      \item Order cost ($S$)
      \item Annual holding cost per unit ($h$)
      \item Unit cost ($c$)
    \end{itemize}
    \item \textbf{Decision variable:} order quantity ($Q$) \hfill \code{\#\# why?}
    \begin{tabular}{lr}
      \begin{minipage}{0.4\textwidth}
        \begin{itemize}
        \item \# Annual order = $\frac{D}{Q}$
        \item Avg Inventory = $\frac{Q}{2}$
        \end{itemize}
      \end{minipage}
    &
      \begin{minipage}{0.55\textwidth} \footnotesize{
      if demand = 240 unit/year, then \\
      \begin{tabular}{||l|rrrrr||} \hline \hline
        qty (Q)    & 1   & 10  & 30 & 40 & 120 \\ \hline
        \# order    & 240 &  24 &  8 & 6  & 2  \\ \hline \hline
      \end{tabular} }
    \end{minipage}\\
    \end{tabular}
    \item[]
    \item \textbf{Total Costs:} material cost + ordering cost + holding cost
  \end{itemize}
  \begin{alignat*}{2}
    TC(Q)  &=\ & \overbrace{c\, D  }^{\mbox{material}}
            + \overbrace{S\,\frac{D}{Q}  }^{\mbox{ordering}}
            + \overbrace{h\,\frac{Q}{2}  }^{\mbox{holding}}
  \end{alignat*}

  \note{
    \noindent\textbf{Purpose:} Definition and derivation of EOQ\par
    \noindent\textbf{Leading Question:} How do the order and average-stock terms change as quantity increases?\par
    \noindent\textbf{Teaching cue:} Compute orders per year and average units for one row before multiplying by their rates.\par
    \begin{itemize}
      \item \textbf{Why Num. Repl} $= \frac{D}{Q}$ based on constant demand
      \item \textbf{Why Avg Inv} $= \frac{Q}{2}$ based on figure min is zero and max is $Q$
      \item example monthly demand is 10 units or 120 units annually.
    \end{itemize}
    \begin{tabular}{lrrrrr} \hline \hline
      case      & C01 & C02 & C03 & C04 & C05 \\ \hline
      order qty & 120 &  30 &  5  &   1 & 9\\
      freq per yr& 1  &   4 &  24 &  240& 13.3  \\ \hline \hline
    \end{tabular}
   }
\end{frame}

%============================================================================================================--
  \begin{frame}
  \frametitle{Holding cost \textbf{VS} Ordering cost} \RemoveLine  %\RemoveHalfLine

  \begin{center}
    \includegraphics[width=1.00\linewidth]{\figureLocation{scm_inv/EOQ_TotalCost.jpg}}\\
  \end{center}
  \begin{flushright}
    \tiny{\textbf{source:} Chase and Jacob. 2011.\hskip10em}\\
  \end{flushright}

  \note{
    \noindent\textbf{Purpose:} Visualize result of EOQ\par
    \noindent\textbf{Leading Question:} Why is the EOQ total-cost curve relatively flat near its minimum?\par
    \noindent\textbf{Teaching cue:} Compare holding and ordering costs on both sides of the minimum.\par
    \begin{itemize}
      \item \textbf{holding cost:} linearly increasing function wrt order qty
      \item \textbf{ordering cost:} non-linearly (inverse) decreasing function wrt order qty
      \item holding cost $=$ ordering cost $\rightarrow$ lowest/minimum total cost
    \end{itemize}
  }
  \end{frame}

%============================================================================================================--
  \begin{frame}
    \frametitle{Derivation of EOQ} \RemoveLine  %\RemoveHalfLine

    \begin{center}
      \begin{block}{\small{Economic Order Quantity}}{}
      \begin{eqnarray*}
        \mbox{Quantity } (Q^{*}) &=& \sqrt{ \frac{2 D\ S }{ h }} \\
        \mbox{Total Cost } \left(TC(Q^{*})\right) &=& c D\ +\sqrt{ 2 D \ S \ h }
      \end{eqnarray*}
      \end{block}
    \end{center}

    \begin{eqnarray*}
    \mbox{where, }   Q        &=& \mbox{Quantity}\\
      TC(Q)    &=& \mbox{Total costs}\\
      D        &=& \mbox{Demand rate}\\
      c        &=& \mbox{Unit purchasing Cost}\\
      S        &=& \mbox{Ordering cost}\\
      h        &=& \mbox{Holding cost}\\
    \end{eqnarray*}

    \note{
    \noindent\textbf{Purpose:} Result of derivation\par
    \noindent\textbf{Leading Question:} How can you verify the EOQ expression without memorizing it?\par
    \noindent\textbf{Teaching cue:} Differentiate relevant annual cost and check the positive stationary point and units.\par
    \begin{itemize}
      \item \textbf{Calculus/FOC} (First Order Condition) $\frac{\partial}{\partial Q} f(Q) = 0$ as slope is equal to zero
    \end{itemize}
    To find $Q^*$
    \begin{alignat*}{4}
      &~~~~& f(Q) &=& c\,D + S\,\frac{D}{Q} + h\,\frac{Q}{2} \\
      \frac{\partial}{\partial Q} f(Q) = 0 &;~~~~&  0 +  S\,(-1)\,\frac{D}{Q^2} + h\,\frac{1}{2} &=& 0 \\
      &~~~~& h\,\frac{1}{2} &=& S\,\frac{D}{Q^2} \\
      &~~~~& Q^2 &=& 2\,S\,D \frac{1}{h}
    \end{alignat*}
    To find $f(Q^*)$
    \begin{alignat*}{2}
      f(Q^*= \sqrt{\frac{2\,S\,D}{h}} ) &=& c\,D + S\,\frac{D}{ \sqrt{\frac{2\,S\,D}{h}} } + h\,\frac{ \sqrt{ \frac{2\,S\,D}{h} } }{2}  \\
           &=& c\,D +  \sqrt{ \frac{S\,D\,h}{2} } + \sqrt{ \frac{S\,D\,h}{2} }
    \end{alignat*}

    }
  \end{frame}

%============================================================================================================--
  \begin{frame}
  \frametitle{Example of EOQ}

    \small{At Allied Electric (HK), The expected annual demand of Samsung TVs is 12,000 units. Administrative, transportation, and receiving form the fixed cost of \$4000 each time an order is places. Each TV cost \$500 to buy and the annual holding cost is 20 \% of value. Find EOQ

    \begin{flushright}
      \tiny{\textbf{source:}Chopra and Meindl. 2010  example 10.1 \hskip10em}\\
    \end{flushright}

    \textbf{Parameters}
    \begin{itemize}
      \item \textbf{Demands (D):} \uncover<2-|handout:0>{12,000 TVs per year}
      \item \textbf{Unit cost (c):} \uncover<2-|handout:0>{500 dollar per unit}
      \item \textbf{Ordering Cost (S):} \uncover<2-|handout:0>{ 4000 dollar per transaction}
      \item \textbf{holding Cost (h):} \uncover<2-|handout:0>{ 500 (0.2)  =  100 dollars per TV per year}
    \end{itemize}

    \uncover<3-|handout:0>{
      \begin{eqnarray*}
        Q^{*} &=& \sqrt{\frac{2 (12000) (4000)}{100}  }\\
              &=& \sqrt{960{,}000} \approx 980
      \end{eqnarray*}
      }
    }

    \note{
    \noindent\textbf{Purpose:} A simple example of EOQ\par
    \noindent\textbf{Leading Question:} What annual holding cost per TV belongs in the formula?\par
    \noindent\textbf{Teaching cue:} Multiply price by 20 percent. Use the corrected radicand 960,000 and round about 979.80 TVs for an operating decision.\par
      \begin{itemize}
        \item \textbf{Demands (D):} 12,000 TVs per year
        \item \textbf{Unit cost (c):} 500 dollar per unit
        \item \textbf{Ordering Cost (S):} 4000 dollar per transaction
        \item \textbf{holding Cost (h):} 500 (0.2)  =  100 dollars per TV per year
        \item \textbf{plug-in:} $Q^{*} = \sqrt{\frac{2 (12000) (4000)}{100} } = \sqrt{960{,}000} \approx 980$
        \item \textbf{solution:} annual ordering and holding = 97,980
      \end{itemize}
    }

  \end{frame}

%============================================================================================================--
  \begin{frame}
  \frametitle{EOQ Insight to inventory control}

  \begin{itemize}
    \item Trade-off between \alert{holding cost} and \alert{ordering cost}
    \begin{minipage}{0.8\textwidth}
      \begin{alertblock}{\small{Key insight}}{}
        Total holding cost = Total ordering cost
      \end{alertblock}
    \end{minipage}
    \item[]
    \item $Q^{*}$ \alert{unaffected} by unit cost of product
    \item $Q^{*}$ can be reduced by
    \begin{minipage}{0.8\textwidth}
      \begin{block}{}{}
      \begin{itemize}
        \item<2-> decreasing \alert{ordering cost} \uncover<3-|handout:0>{ $\rightarrow$ good idea}
        \item<2-> increasing handling cost \uncover<3-|handout:0>{ $\rightarrow$ bad idea}
      \end{itemize}
      \end{block}
    \end{minipage}
    \item[]


    \item Using EOQ creates \alert{bullwhip effect}
  \end{itemize}

  \NextLine
  \textbf{Limitation of EOQ:} rank from possible to impossible
  \begin{itemize}
    \item<3-> instant replenishment
    \item<3-> nearly constant demand
    \item<3-> zero and constant lead time
  \end{itemize}

    \note{
    \noindent\textbf{Purpose:} Key result from EOQ\par
    \noindent\textbf{Leading Question:} When can a nearby feasible order size be preferable to the formula quantity?\par
    \noindent\textbf{Teaching cue:} Compare annual relevant costs after packaging, truckload, or minimum-order constraints.\par
    \begin{itemize}
      \item \textbf{Property} quality determined by EOQ is robust
      \item \textbf{Key Question:} why holding cost $=$ ordering cost $\rightarrow^{trade-off}$ idea of many heuristics
      \item \textbf{Result} how to reduce ordering size; Toyota's setup time
      \item \textbf{Limitation}
    \end{itemize}


    }
  \end{frame}

%============================================================================================================--
  \begin{frame}
  \frametitle{Academic Extension of EOQ Model  \forIE}

  \textsc{2104611: Inventory Analysis}
  \begin{description}
    \item[Production] Replenishment is \alert{not instantaneous}
    \item[Multiple Products] Combine dependent products to share ordering cost
    \item[Multi-echelon] Upstream EOQ is \alert{multiple of downstream EOQ}
    \item[Safety stock] Additional quantity to protect \alert{surging in demands}
  \end{description}

  \NextLine
  \textsc{Next Example}
  \begin{description}
    \item[Quantity Discount] Buy more to reduce total costs
    \item[Storage Constraint] Warehouse has limited \alert{storage space}
  \end{description}

  %%%-----
   \note{
    \noindent\textbf{Purpose:} Expansion of EOQ, academically (study in Inventory Analysis)\par
    \noindent\textbf{Leading Question:} Which assumption breaks when the supplier offers a price break?\par
    \noindent\textbf{Teaching cue:} Include acquisition cost and evaluate feasible price tiers and breakpoint candidates.\par
    \begin{itemize}
      \item \textbf{Production} replenish can be constantly done; require rate greater than demand, e.g., batch production
      \item \textbf{Multiple Products}use power-of-two (order in sequence of 2) to cope with order depending SKU demands
      \item \textbf{Quantity Discount} \& \textbf{Storage Capacity}think and understand beyond derivation $\rightarrow$ see example
      \item \textbf{Multi -echelon}combine qty and synchronize upstream-downstream inventory; think and analyze the whole SCM
      \item \textbf{Safety stock}assume demand is normally distributed and convert probability into service level
    \end{itemize}

  }
  \end{frame}

%============================================================================================================--
  \begin{frame}
  \frametitle{Example: Ordering carton}

  \small{A bottler company must order 30,000 corrugated boxes from a cardboard suppler, annually. Currently, the suppler sells cardboard at 10 THB per box and charges a delivery cost at 400 THB per transaction.
  If the company estimate its internal handling costs at 15\% of its values. Answer the following questions

  \begin{itemize}
    \item[a] What is the EOQ and total logistics costs?
    \item[b] If the company has only 2000 boxes storage space within it compound the factory, how is the solution changed?
    \item[c] If the company stores additional boxes at an outside warehouse at cost 1 THB per box per year, how is the solution changed?
    \item[d] If the suppler waives the delivery charge when the company orders more than 30,000 THB per order, how is the solution changed?
  \end{itemize}
  }

  \note{
    \noindent\textbf{Purpose:} Example of Adv EOQ question\par
    \noindent\textbf{Leading Question:} How should the carton answer change when the in-house limit starts to bind?\par
    \noindent\textbf{Teaching cue:} Compare the unconstrained quantity with the storage limit before evaluating external capacity.\par
    \begin{itemize}
      \item $<<$Explain Questions$>>$
      \item problem is sequential. Question D includes all info of Questions A, B, and C
      \item b with a $\rightarrow$ c with (a,b)  $\rightarrow$ d with (a,b,c)
    \end{itemize}
  }
  \end{frame}

%============================================================================================================--
  \begin{frame}
  \frametitle{Solution: Ordering carton}

  \begin{itemize}
    \item[] \textsc{Parameters}
    \begin{itemize}
      \item \textbf{Demands (D):} \uncover<2-|handout:0>{30,000 boxes per year}
      \item \textbf{Ordering Cost (S):} \uncover<2-|handout:0>{ 400 THB per transaction}
      \item \textbf{holding Cost (h):} \uncover<2-|handout:0>{ 10 (0.15)  =  1.5 THB per box per year}
    \end{itemize}
    \item[] \textsc{Solution}
    \begin{itemize}
    \item[a)]<2-> \textbf{EOQ ($Q^*$): }  \uncover<2-|handout:0>{ $\sqrt{\frac{2\,30,000\,400}{1.5}}$  =  4,000 boxes at cost 6,000}
    \item[b)]<2-> \textbf{Limited Space: } \uncover<2-|handout:0>{$TC(Q=2,000) =  \frac{30,000}{2000} (400)+ \frac{1}{2} 2000 \, (1.5) = 7,500$}
    \item[c)]<3-> \textbf{Rent: } \uncover<2-|handout:0>{If renting at $\hat{h}$ THB per box-year (permanent/ no transferring)
    \[  TC(Q) = S \frac{D}{Q} + (Q-2000)\,\hat{h} + h \, \frac{1}{2}Q  \]   }
    \item<3-|handout:0>[ ]  \hskip5em $\frac{DS}{Q^2} = \frac{1}{2} h + \hat{h}$; \hskip3em $TC(Q^* = 2619)= 7,165$ ~~~
    \item[d)]<3-> \textbf{Free shipping:} \uncover<2-|handout:0>{$TC(3{,}001) = (3{,}001-2{,}000)(1) + \frac{1}{2}(1.5)(3{,}001) = 3{,}251.75$ THB/year}
    \end{itemize}
  \end{itemize}


  \note{
    \noindent\textbf{Purpose:} How to answer each question\par
    \noindent\textbf{Leading Question:} Which carton order size first meets a strictly greater-than free-delivery threshold?\par
    \noindent\textbf{Teaching cue:} At 10 THB per box the first eligible whole order is 3,001. Compare the stated rental scenario's 3,251.75 THB annual cost with paid-delivery alternatives.\par
    \begin{itemize}
    \item[a)] \textbf{EOQ ($Q^*$): }  $\sqrt{\frac{2\,30,000\,400}{1.5}}$  =  4,000 boxes at cost 6,000
    \item[b)] \textbf{Limited Space: } $TC(Q=2,000) =  \frac{30,000}{2000} (400)+ \frac{1}{2} 2000 \, (1.5) = 7,500$
    \item[c)] \textbf{Rent: } If renting at $\hat{h}$ THB per box-year (permanent/ no transferring)
    \[  TC(Q) = S \frac{D}{Q} + (Q-2000)\,\hat{h} + h \, \frac{1}{2}Q  \]
    \item[ ]  \hskip5em $\frac{DS}{Q^2} = \frac{1}{2} h + \hat{h}$; \hskip3em $TC(Q^* = 2619)= 7,165$ ~~~
    \item[d)] \textbf{Free shipping:} $TC(3{,}001) = (3{,}001-2{,}000)(1) + \frac{1}{2}(1.5)(3{,}001) = 3{,}251.75$ THB/year
    \end{itemize}
  }
  \end{frame}

%============================================================================================================--
  \begin{frame}
  \frametitle{Can EOQ apply to real world? }

  \begin{itemize}
    \item \textbf{Demand:} stochastic $\rightarrow$ known distribution
    \item \textbf{Lead time:} constant(qty)/random
    \item \textbf{How:} continuous review of current inventory
    \item \textbf{When:} enough for \alert{expected sale} till get a new order
    \item \textbf{Safety Stock:} inventory for randomness till waiting a new order
  \end{itemize}

  \NextLine
  \begin{alertblock}{\small{\textsc{Example of Policy}}  \forIE}{}
    \begin{itemize}
      \item[]  $(Q,q)$ policy is used for order
      \begin{itemize}
        \item q: reorder point, expected demand during lead time + safety stock
        \item Q: order quantity derived from EQO model (using expected sale) + safety stock
      \end{itemize}
      \item $\exists$ chance of out-of-stock $\rightarrow$ \alert{service level}
    \end{itemize}
  \end{alertblock}

  \note{
    \noindent\textbf{Purpose:} Problem of EOQ in practice and how practically overcome them\par
    \noindent\textbf{Leading Question:} Which practical constraint could move the recommended lot size away from EOQ?\par
    \noindent\textbf{Teaching cue:} Test a supplier minimum or delivery schedule and compare feasible annual costs.\par
    \begin{itemize}
      \item[] $<<$Explain Bullets$>>$
      \item \textbf{Leadtime:} in many case leadtime depends on ordering qty
      \item \textbf{Feeback:} use inventory qty as info for when to order
      \item \textbf{Safety Stock:} use excess inventory as buffer against uncertainty during ordering
    \end{itemize}
  }
  \end{frame}

%============================================================================================================--
  \begin{frame}
  \frametitle{Real world: Random demand and lead time}

  \begin{center}
    \includegraphics[width=1.00\linewidth]{\figureLocation{scm_inv/RealWorldInventory.jpg}}\\
  \end{center}
  \begin{flushright}
    \tiny{\textbf{source:} Chase and Jacob. 2011.\hskip10em}\\
  \end{flushright}

  \note{
    \noindent\textbf{Purpose:} inventory curve of real world\par
    \noindent\textbf{Leading Question:} Which uncertainty contributes to demand during lead time?\par
    \noindent\textbf{Teaching cue:} Separate demand-rate variation from lead-time variation before selecting a safety-stock method.\par
    \begin{itemize}
      \item \textbf{Case01:} no consume safety stock $\rightarrow$ prob 0.50 if demand is normal.
      \item \textbf{Case02:} consume safety stock, but remain $\rightarrow$ purpose of safety stock (depend on service level)
      \item \textbf{Case03:} insufficient; too expensive to support all demand
      \item \textbf{Question:} What is a review method?
    \end{itemize}
  }
  \end{frame}

%============================================================================================================--
  \begin{frame}
  \frametitle{Safety stock}

  \begin{itemize}
    \item \textbf{Idea:} demand is unpredictable; if $\exists \!\!\!/$ buffer inventory, lost sale
    \item \textbf{What:} buffer inventory for \alert{unexpected demands} during ordering
    \item \textbf{Factors} \alert{variability} of demands and lead time
    \item \textbf{Formulation:} safety stock = \alert{service level} $\times$ \alert{variation}
    \[ \mbox{avg. demand$_{LT}$} =  \bar{D} \,\bar{L} \quad ; \quad \mbox{var. demand$_{LT}$} = \bar{L} \,\sigma_D + \bar{D}^2 \,\sigma_L \]
    \item \textbf{De-mystify:} safety stock $\neq$ demands during lead time
    \item \textbf{Reduce safety stock:} priority to reduce safety stock
    \begin{itemize}
      \item reduce \alert{variability} of lead time
      \item reduce \alert{expected} lead time
      \item reduce \alert{variability} of demand
    \end{itemize}
  \end{itemize}


  \note{
    \noindent\textbf{Purpose:} Introduce of Safety Stock\par
    \noindent\textbf{Leading Question:} What service measure is the buffer intended to protect?\par
    \noindent\textbf{Teaching cue:} Distinguish cycle service from fill rate and connect buffer size to the distribution.\par
    \begin{itemize}
      \item \textbf{Formula} base on independence of LT and demand distributions; (one of assignment in Conditional Probability Class)
      \begin{itemize} \scriptsize{
        \item \textbf{expect total qty} is product of avg LT and demand; property of expected value
        \item \textbf{expect total variance} is more complex and require conditional probability to show. however, $\sigma_L$ is the first priority reduce }
      \end{itemize}
      \item \textbf{no LT, no safety stock:} e.g., water in your home
    \end{itemize}
  }
  \end{frame}

%============================================================================================================--
\section[Review System]{Inventory Review System}
%============================================================================================================--
  \begin{frame}
  \frametitle{How often inventory should be determined?}

  \textsc{Periodic review system} \hfill \hyperlink{periodic}{Pg: Periodic Review}
  \begin{itemize}
    \item \textbf{Idea:} \alert{checking inventory} at regular period
    \item \textbf{Example:} end-of-fiscal year
    \item \textbf{Pro:} no tool, less effort to review
    \item \textbf{Evidence:} order duration
  \end{itemize}

  \NextLine
  \textsc{Continuous review system}  \hfill \hyperlink{cont}{Pg: Continuous Review}
  \begin{itemize}
    \item \textbf{Idea:} continuously \alert{tracking inventory}
    \item \textbf{Example:} two-bin system
    \item \textbf{Pro:} less avg inventory
    \item \textbf{Evidence:} order date, trigger qty level
    \item \textbf{Note:} Tool=IT to improve speed and \alert{accuracy}
  \end{itemize}

  %%%------
  \note{
    \noindent\textbf{Purpose:} Introduce inventory Review System\par
    \noindent\textbf{Leading Question:} How does review frequency change the period inventory must protect?\par
    \noindent\textbf{Teaching cue:} Compare the protection interval for periodic and continuous review.\par
    \begin{itemize}
      \item \textbf{Continuous:} instant tracking of inventory data
      \begin{itemize} \scriptsize{
        \item \textbf{Scope of time:} realtime $\rightarrow$ when to order $\rightarrow$ day
        \item \textbf{Practice:} IT controls over incoming and outgoing sku; kanban of GR in hospital; 2-bin system in rest room }
      \end{itemize}
      \item \textbf{Periodic:} offline checking of inventory availability
      \begin{itemize} \scriptsize{
        \item \textbf{Normal Condition:} few suppliers, limited purchaser, inaccurate inventory qty and location
        \item \textbf{Practice:} ask warehouse staff to check inventory before order}
      \end{itemize}
    \end{itemize}
  }
  \end{frame}

%============================================================================================================--
  \begin{frame}[label=periodic]
  \frametitle{Periodic review \forIE}

  \begin{center}
    \includegraphics[width=0.90\textwidth]{\figureLocation{scm_inv/fixedInterval.png}}\\
  \end{center}

    \note{
    \noindent\textbf{Purpose:} what periodic review look like\par
    \noindent\textbf{Leading Question:} How do you calculate the order at a scheduled review?\par
    \noindent\textbf{Teaching cue:} Subtract inventory position from the target and check on-order units and backorders.\par
    \begin{itemize}
      \item \textbf{inter-order:} almost constant or multiple of one another; occurring on certain day of wekk
      \item \textbf{trigger qty:} higher gap
      \item \textbf{note:} never use incoming qty for indicator as it relates with shipment
    \end{itemize}
  }
  \end{frame}

%============================================================================================================--
  \begin{frame}[label=cont]
  \frametitle{Continuous review  \forIE}

  \begin{center}
    \includegraphics[width=0.90\textwidth]{\figureLocation{scm_inv/reOrderPoint.png}}\\
  \end{center}

    \note{
    \noindent\textbf{Purpose:} what periodic review look like?\par
    \noindent\textbf{Leading Question:} What changes the inventory position and triggers a replenishment in continuous review?\par
    \noindent\textbf{Teaching cue:} Track on-hand, on-order, and backorders through one event.\par
    \begin{itemize}
      \item \textbf{inter-order:} inconsistant
      \item \textbf{trigger qty:} almost constant after pass certain qty
      \item \textbf{note:} never use incoming qty for indicator as it relates with shipment
    \end{itemize}
  }
  \end{frame}

 %============================================================================================================--
  \begin{frame}
  \frametitle{Do we have good inventory management?}

  \textsc{Symptoms of bad inventory control}
  \begin{itemize}
    \item insufficient inventory
    \item non-moving/ dead stock/ obsolete
    \item inaccurate stock
  \end{itemize}

  \NextLine
  \textsc{KPI}
  \begin{itemize}
    \item \textbf{Fill rate:} $\frac{\mbox{product shipped}}{\mbox{product order}}$
    \item \textbf{Inventory turnover:} $\frac{\mbox{annual sale}}{\mbox{inventory}}$
    \item \textbf{Inventory accuracy:} $1 - \frac{|\mbox{actual} - \mbox{account} | }{\mbox{actual}}$
  \end{itemize}

    \note{
    \noindent\textbf{Purpose:} Symptom and KPI\par
    \noindent\textbf{Leading Question:} How can a lower inventory balance conceal worse service?\par
    \noindent\textbf{Teaching cue:} Pair a stock measure with fill rate, shortage, or backorder duration.\par
    \begin{itemize}
      \item \textbf{Dead Stock:} most simple problem, but require financial wills
      \item \textbf{Inaccurate:} required consistency to solve
      \item \textbf{Insufficient:} most problematic because need some other alternative to avoid risks
      \item \textbf{Alternative:} borrow, back-back contract, delay/partial shipment/backlog, common items change form
      \item[] \textbf{\textsc{KPI:} Key performance index}
      \item \textbf{All:} all KPI has many different levels: order/sku/qty
      \item \textbf{Fill rate:} measuring available SKUs
      \item \textbf{Inventory turnover/day:} measuring sales/inventory
      \item \textbf{Inventory accuracy:} data in system VS real inventory
    \end{itemize}
  }
  \end{frame}

 %============================================================================================================--
  \begin{frame}
  \frametitle{Inventory turnover}

  \begin{center}\small{
    \begin{tabular}{lrrrrr}\hline
      	                  & 2011	& 2010	& 2009	& 2008	& 2007 \\ \hline
      \textbf{Kohl's} \\
      Revenue	            & 18391	& 17178	& 16389	& 16474	& 15544 \\
      Cost of Goods Sold	& 11359	& 10679	& 10332	& 10459	& 9890\\
      Inventory	          & 3036	& 2923	& 2799	& 2856	& 2588 \\
      Net Income	        & 1114	& 991	& 885	& 1084	& 1109 \\ \hline
      	& 	& 	& 	& 	& \\
      \textbf{Walmart} \\
      Revenue	            & 418952	& 405046	& 401244	& 374526	& 344992\\
      Cost of Goods Sold	& 307646	& 297500	& 299419	& 280198	& 258693\\
      Inventory	          & 36318	  & 33160	  & 34511	& 35180	& 33685 \\
      Net Income	        & 16389	  & 14335	  & 13188	& 12884	& 12036 \\ \hline
    \end{tabular} }
  \end{center}

  \[
  \mbox{Inventory turnover} = \frac{11359}{3036} = 3.74
  \]
    \note{
    \noindent\textbf{Purpose:} example of inventory analysis in AR\par
    \noindent\textbf{Leading Question:} Which units and time period define inventory turnover?\par
    \noindent\textbf{Teaching cue:} Use consistent valuation for cost of goods sold and average stock before comparing firms.\par
    \begin{itemize}
      \item in addition of inventory control, we can apply the analysis in annual report and comparing company
    \end{itemize}
  }
  \end{frame}

%============================================================================================================--
\section[NewsVendor]{Stochastic Demand with Newsvendor Model}
%============================================================================================================--

  \begin{frame}
  \frametitle{Perishable Products}

  \begin{itemize}
    \item \textbf{Product:} good short period and single order
    \item \textbf{Example:} newspaper, magazines, flowers, produces, seasonal clothing
  \end{itemize}

  \NextHalfLine
  \textbf{Example}\\
  \small{ A newspaper costs ($c$) \$1 and can be sold for ($p$) \$2 with savage value ($v$) zero. What is the profit ($\pi(q)$). if we order ($q$) \alert{20} today and demand ($d$) is}:
  \begin{center}
  \begin{tabular}{lrrrr} \hline
    Demand ($d$) & Sold & Leftover & Lost sale & Profit ($\pi$) \\ \hline
      20         & 20   & 0        & 0         & $ 2 (20) - 1 (20) + 0(0) = \$20$ \\
      15         & 15   & 5        & 0         & $ 2 (15) - 1 (20) + 0(5) = \$10$ \\
      30         & 20   & 0        & 10        & $ 2 (20) - 1 (20) + 0(0) = \$20$ \\ \hline
  \end{tabular}
  \end{center}

  \begin{alertblock}{}{}
    \begin{eqnarray*}
      \pi(q)\  = \overbrace{ p \times \min(d,q)}^{\mbox{\tiny{revenue}}} - \overbrace{ c \times  q}^{\mbox{\tiny{purc.cost}}} + \overbrace{v \times  max(0,q-d)}^{\mbox{\tiny{savage value} }}
    \end{eqnarray*}
  \end{alertblock}

  \note{
    \noindent\textbf{Purpose:} Introduce of inventory concept for perishable items\par
    \noindent\textbf{Leading Question:} What makes an unsold perishable unit costly at the end of the selling period?\par
    \noindent\textbf{Teaching cue:} Identify salvage, disposal, and holding effects separately from acquisition cost.\par
    \begin{itemize}
      \item \textbf{Story:} A 30-house village that has only single newspaper shop. In each day, a shop owner need to order newspaper at cost of 1.00USD to sell to villagers at price 2.00USD. If there is left over, he lose 1.00USD money as newspaper cannot be sold the other day nor returned.   Suppose that he decided to order 20 newspaper, there are three possible cases:
      \item \textbf{Case01: Luckily:} sold-out with 20USD profit. incur no lose and no further profit can be made.
      \item \textbf{Case02: Too much:} 5 newspaper left-over with 10USD profit. incur 5.0USD lose, can made more profit by order 1 newspaper less down to the demand.
      \item \textbf{Case03: Too few:} sold-out with 20USD profit similar to Case01. However, he can made more profit by order 1 newspaper more up to the demand.
      \item \textbf{Observation 1:} meet demand $\rightarrow$ max profit. However, demand is uncertainty; hence, pattern/distribution of demand matter
      \item \textbf{Observation 2:} profit margin $\propto$  risk as if cost is low or price is high, more profit margin. $\rightarrow$ more risk.
      \item \textbf{Observation 3:} risk $\propto$ distribution of demand
      \item \textbf{Observation 0:} if salvage $<!\!\!\!\!\/$ cost $<!\!\!\!\!\/$ price, this setting does not work as owner can arbitrage.
      \item \textbf{Equation:} what is the relationship look like, $\min(d,q)$ and $\max(0,q-d)$ are available newspaper and left-over qty.
    \end{itemize}
  }
  \end{frame}

%============================================================================================================--
  \begin{frame}
  \frametitle{Optimal Newsvendor}

    \footnotesize{
    \begin{center}
      \begin{block}{Newsvendor}{}
      \begin{eqnarray*}
        \mbox{Expected total Profit}  &=& \mbox{Expected net revenue} - \mbox{Total net investment}  \\
        \pi(q) &=& (p-v) \cdot S(q) +v \cdot \mu(q)  - c\cdot q  \\
        \mbox{Quantity } (q^{*}) &=& F^{-1}\left(\frac{p - c}{p - v} \right)  \hskip2ex =  F^{-1}\left(\frac{p - c}{p \alert{-c+c} -v } \right) \\
                                 &=& F^{-1}\left( \frac{C_{u} }{C_u + C_o } \right)
      \end{eqnarray*}
      \end{block}
    \end{center}
    }
    \RemoveLine
    \footnotesize{
    \uncover<2->{

    \begin{eqnarray*}
    \mbox{where, }   \pi(q)    &=& \mbox{Total profit}\\
      F(\cdot)  &=& \mbox{Cumulative probability function}\\
      S(q)      &=& \mbox{Expected units sold}, S(q) = q - \int_{0}^{q}F(y)dy \\
      \mu(q)    &=& \mbox{Expected units unsold}, \mu(q) = \int_{0}^{q}F(y)dy\\
      C_u       &=& \mbox{Cost of under stock}, C_u = p-c\\
      C_o       &=& \mbox{Cost of overstock}, C_o = c-v
    \end{eqnarray*}
    }}


  \note{
    \noindent\textbf{Purpose:} Introduce derivation of newsvendor problem\par
    \noindent\textbf{Leading Question:} Why is the newsvendor critical fractile a service probability?\par
    \noindent\textbf{Teaching cue:} Compare marginal underage and overage loss and read the quantile of the demand distribution.\par
    \begin{itemize}
      \item \textbf{Explain:} base on a careful Math, the cdf of demand function relates to profit margin
      \item \textbf{Critical Ration:} ($\frac{p - c}{p - v}$). If $v=0$, this term is profit margin!
      \item \textbf{Expected unsold:} $\mu(q) = \int_0^{q} F(y) dy$
      \item \textbf{Expected sold} $S(q) = q - \int_0^{q} F(y) dy$
      \item \textbf{Tech Opt. Note:} need to write \alert{equation} as expected expression and argue that limit is finite so FOC is finite. Hence, swap between integrate and differentiate operations
    \end{itemize}
  }

  \end{frame}
%============================================================================================================--

\begin{frame}
  \frametitle{Example of Newsvendor}

  \small{Sportmart needs to determine quantities of skis for the winter. Based on historical demand, the demand of ski is normally distributed with mean ($\mu$) 350 and sd ($\sigma$) 100. Each pair of ski costs \$100 and retail for \$250. The unsold skis are disposed for \$85 assuming it costs \$5 to hold inventory for the entire winter.
   }
   \begin{flushright}
    \tiny{\textbf{source:}Chopra and Meindl. 2010  example 12.1 \hskip10em}\\
  \end{flushright}

  \NextHalfLine
  \textbf{Steps}
   \begin{itemize}
    \item[0] \textbf{Model:} Newsvendor model $\mathcal{N}(\mu=350,\sigma=100)$
    \item[1] \textbf{Parameter:} $p = \$250, \ c = \$100, \ v = 85-5 = \$80$
    \item[] \hskip3em $C_{u} = 250-100 = \$150$ and $C_o = 100 - 80 = \$20$
    \item[2] \textbf{Critical fractile:} $F(q^*) = \frac{150}{150 + 20} \approx 0.88235$
    \item[3] \textbf{Quantity:} $q^* = 350 + 100\Phi^{-1}(150/170) \approx 468.68$ pairs
  \end{itemize}

  \note{
    \noindent\textbf{Purpose:} Example of newsvendor question\par
    \noindent\textbf{Leading Question:} Where does rounding the critical ratio change the ski order quantity?\par
    \noindent\textbf{Teaching cue:} Use 150 divided by 170 before converting to a normal quantile. The continuous optimum is about 468.68 pairs.\par
    \begin{itemize}
      \item \textbf{cost of overstock:} $(C_o=c-v)$ is the net loss on one excess unit
      \item \textbf{cost of understock:} $(C_u=p-c)$ is the contribution margin lost on one unfilled unit
      \item \textbf{cdf of normal distribution} $\Phi^{-1}(150/170)\approx1.187$. Compare adjacent feasible integers if orders must be whole pairs.
    \end{itemize}
  }
\end{frame}

%============================================================================================================--
  \begin{frame}
  \frametitle{Insight to Inventory Control}

  \begin{itemize}
    \item Trade-off between \alert{opportunity cost} and \alert{overstocking cost}
    \item $q^*$ is equilibrium overage and underage costs are balanced
    \item Assuming $p > c > v$
    \item $\frac{C_u}{C_u +C_o}$ is between $[0,1]$ similar to probability
    \item Large $C_u$ comparing to $C_o$ $\rightarrow$ \alert{huge purchase}
    \item Small $C_u$ comparing to $C_o$ $\rightarrow$ \alert{little purchase}
  \end{itemize}



  \note{
    \noindent\textbf{Purpose:} Key concept of newsvendor\par
    \noindent\textbf{Leading Question:} Which policy fits a stable recurring item versus a one-season item?\par
    \noindent\textbf{Teaching cue:} Use the time horizon, flexibility to reorder, and salvage assumptions to choose the policy.\par
      \begin{itemize}
        \item \textbf{trade-off:} $q^*$ is equilibrium of the trade-off between \alert{opportunity cost} and \alert{overstocking cost}
        \item \textbf{Range:} if an owner makes no profit ($c=p$), he takes no risk. and if owner bears no loss ($c=v$), he takes all risk
        \item \textbf{arbitrage:} what happen when salvage $>$ cost?
        \item \textbf{Huge purchase/All risk} $1 \approx \frac{C_u}{C_u + C_o} = \frac{p-c}{p-v} \rightarrow (p-c) \approx (p-v) \rightarrow c \approx v $ or $c \approx 0 \& v = 0$
        \item \textbf{Few purchase/No risk} $0 \approx \frac{C_u}{C_u + C_o} = \frac{p-c}{p-v} \rightarrow p=c$
      \end{itemize}

  }
  \end{frame}



%================================================================================================-

\begin{frame}
  \frametitle{Reference}

  \begin{footnotesize}
    \bibliography{\refLocation{Coordination}}
  \end{footnotesize}

  \note{
    \noindent\textbf{Purpose:} Show reference\par
    \noindent\textbf{Leading Question:} What would you check before reusing a figure, fact, or model from this lecture?\par
    \noindent\textbf{Teaching cue:} Trace a citation to its original source and verify the year, scope, and assumptions.\par
    \begin{itemize}
      \item Confirm the source and edition before assigning a reading.
    \end{itemize}
  }

\end{frame}

%==============================================================================================================
%\section[Inventory Game]{Pharmaceutical Inventory Game}
%==============================================================================================================
\begin{frame}
  \frametitle{How to play Pharmaceutical Inventory Game}

  \begin{center}
    \only<1 |handout:0>{ \includegraphics[width=0.80\textwidth]{\figureLocation{inv_game/gameRule.png}} }
    \only<2 |handout:1>{ \includegraphics[width=0.80\textwidth]{\figureLocation{inv_game/gamePlay.png}} }
  \end{center}
  \begin{flushright}
  \begin{tabular}{cl}
    \includegraphics[width=0.10\textwidth]{\figureLocation{inv_game/ChulaInvGame_QR.png}}  &
    \href{https://tinyurl.com/ChulaInvGame}{\code{https://tinyurl.com/ChulaInvGame}} \\
  \end{tabular}
  \end{flushright}


  \note{
    \noindent\textbf{Purpose:} Alternative before inventory storage\par
    \noindent\textbf{Leading Question:} Which decision does the inventory game ask the team to make each period?\par
    \noindent\textbf{Teaching cue:} Confirm the decision order, feedback timing, and how costs and service are recorded.\par
    \begin{itemize}
      \item \textbf{Str: Potent. Altern.:} what consumption? e.g., media $\rightarrow$ internet steaming
      \item \textbf{Str: Design. Product:} design to reduce component? family group, common VS module; decided by mrkt/designer
      \item \textbf{Deploy: Which items?} risk VS space \& money; decided by
      \item \textbf{Deploy: What form?} common for every model VS module for quick response
      \item \textbf{Deploy: Where/Max Qty} money VS responce ;
      \item \textbf{Replenish: Check/Qty/When:} common oper. decision; decided by
    \end{itemize}
  }
\end{frame}

%==============================================================================================================
\begin{frame}
  \frametitle{Historical Data}

  \begin{center}
    \only<1 |handout:0>{
    \begin{tabular}{cc}
      \includegraphics[width=0.45\textwidth]{\figureLocation{inv_game/hist_BW_data.png}}
      &
      \includegraphics[width=0.45\textwidth]{\figureLocation{inv_game/hist_DS_data.png}} \\
    \end{tabular}  }

    \only<2 |handout:0>{ \includegraphics[width=0.80\textwidth]{\figureLocation{inv_game/hist_TS_data.png}} }
  \end{center}

  \note{
    \noindent\textbf{Purpose:} Alternative before inventory storage\par
    \noindent\textbf{Leading Question:} What portion of the historical demand should be held out for a policy check?\par
    \noindent\textbf{Teaching cue:} Separate estimation history from later evaluation and retain the demand units.\par
    \begin{itemize}
      \item \textbf{Str: Potent. Altern.:} what consumption? e.g., media $\rightarrow$ internet steaming
      \item \textbf{Str: Design. Product:} design to reduce component? family group, common VS module; decided by mrkt/designer
      \item \textbf{Deploy: Which items?} risk VS space \& money; decided by
      \item \textbf{Deploy: What form?} common for every model VS module for quick response
      \item \textbf{Deploy: Where/Max Qty} money VS responce ;
      \item \textbf{Replenish: Check/Qty/When:} common oper. decision; decided by
    \end{itemize}
  }
\end{frame}

%==============================================================================================================
%\section[Take Away]{What do we learn from the game?}
%============================================================================================================
\usetikzlibrary{arrows,positioning,decorations.pathmorphing,calc}
\renewcommand{\textFont}[1]{\footnotesize{#1}}
\tikzstyle{lvBlock}    =[rectangle,draw=black!80,thick,text centered,text width=25mm,text height=5mm] %rounded  corners=3mm,
\tikzstyle{detailBlock}=[rectangle,draw=black!80,text width=70mm,text height=5mm]
%=================================================================================================
\begin{frame}
  \frametitle{What are products in the game?}

  \uncover<1 |handout: 0>{
  \begin{center}
    \begin{tabular}{cc}\\
        \includegraphics[height=0.30\textheight]{\figureLocation{inv_game/prodA.png}}
      &
        \includegraphics[height=0.30\textheight]{\figureLocation{inv_game/prodB.png}}
      \\
        \tiny{\textbf{A: 45GWs}: White Strawberry Extract }
      &
        \tiny{\textbf{B: 30MLBkc}: Pure Natural Aromatherapy Skin  }
      \\ %%%%%
        \includegraphics[height=0.30\textheight]{\figureLocation{inv_game/prodC.png}}
      &
        \includegraphics[height=0.30\textheight]{\figureLocation{inv_game/prodD.png}}
      \\
        \tiny{\textbf{C: DGLb530G}: Guava Leaf Extract Juice }
      &
        \tiny{\textbf{D: CL515G}: Calendula extract liquid    }
      \\ %%%%%
    \end{tabular}
  \end{center}
  }


  \note{
    \noindent\textbf{Purpose:} Alternative before inventory storage\par
    \noindent\textbf{Leading Question:} Why might the game require a different policy for each product?\par
    \noindent\textbf{Teaching cue:} Compare variability, lead time, unit value, and service consequences without assuming equal parameters.\par
    \begin{itemize}
      \item \textbf{Str: Potent. Altern.:} what consumption? e.g., media $\rightarrow$ internet steaming
      \item \textbf{Str: Design. Product:} design to reduce component? family group, common VS module; decided by mrkt/designer
      \item \textbf{Deploy: Which items?} risk VS space \& money; decided by
      \item \textbf{Deploy: What form?} common for every model VS module for quick response
      \item \textbf{Deploy: Where/Max Qty} money VS responce ;
      \item \textbf{Replenish: Check/Qty/When:} common oper. decision; decided by
    \end{itemize}
  }
\end{frame}

%=================================================================================================
\begin{frame}
  \frametitle{Important Insights: Cost VS Service }

  \begin{itemize}
    \item \textbf{Focus on what matter:} many not important; only 20\% matter
    \item \textbf{Opportunity Cost:} inventory cost = opportunity, not real money
    \item \textbf{Information = Power:} demand forecasting, parameter accuracy
    \item \textbf{Product Matter:} know products, nature/relationship, regulation (GDP)
    \item \textbf{Move to next phrase:} not just money, but customer, internal growth
  \end{itemize}

  \NextLine
  \textsc{Quick-win in inventory management:}
  \begin{itemize}
    \item<2-> \textbf{Purchasing Power:} use IT or contracts, e.g, VMI, Collaborative Purchasing (CP), paid-per-used
    \item<2-> \textbf{Sharing \& Borrowing:} build alliance, pool risk with \alert{visibility}
    \item<2-> \textbf{Right Dispensary Unit:}  blister packs, box, bottle
  \end{itemize}


  \note{
    \noindent\textbf{Purpose:} Alternative before inventory storage\par
    \noindent\textbf{Leading Question:} Which inventory actions improved service and which only shifted cost?\par
    \noindent\textbf{Teaching cue:} Have teams explain one measured trade-off and one change they would test next round.\par
    \begin{itemize}
      \item \textbf{Str: Potent. Altern.:} what consumption? e.g., media $\rightarrow$ internet steaming
      \item \textbf{Str: Design. Product:} design to reduce component? family group, common VS module; decided by mrkt/designer
      \item \textbf{Deploy: Which items?} risk VS space \& money; decided by
      \item \textbf{Deploy: What form?} common for every model VS module for quick response
      \item \textbf{Deploy: Where/Max Qty} money VS responce ;
      \item \textbf{Replenish: Check/Qty/When:} common oper. decision; decided by
    \end{itemize}
  }
\end{frame}

%================================================================================================-
\end{document}






%============================================================================================================--
  \begin{frame}
  \frametitle{Extension and Limitation}

  \textbf{Extension of Newsvendor Model}
  \begin{description}
    \item[Quantity Discount] Computing multiple newsvendor and selecting highest profit
    \item[Designing Penalty] Given service level, find $C_u$
  \end{description}

  \NextLine
  \textbf{Limitation of Newsvendor Model}
  \begin{itemize}
    \item \alert{demand distribution} and/or historical data
  \end{itemize}

    \note{
      \begin{itemize}
        \item why holding cost $=$ ordering cost
        \item how to reduce ordering size; Toyota's setup time
      \end{itemize}

    }
  \end{frame}
